\section{Security Analysis}
\label{sec:security-analysis}

\subsection{Why Representation Granularity Matters}

For execution context $u$, let $\mu(u)$ be its committed effects and let
$\gamma(u)$ contain policy-relevant metadata such as effect provenance.  The
ideal commit view is $\omega(u)=(\mu(u),\gamma(u))$.  Each policy context
$P\in\mathcal P$ induces an ideal decision $I_P(\omega(u))$.  For example, our
finite effect-set experiments use $I_B(\omega(u))=[\mu(u)\preceq B]$; an ACL,
delegation state, or quota snapshot could parameterize another predicate.  The
policy snapshot used for the check must also govern the commit.  A monitor sees
$\rho(u)$ rather than $\omega(u)$; a descriptor instantiates
$\rho_d(u)=A_{d_t}(u)$.

\begin{definition}[Authorization sufficiency]
\label{def:authorization-equivalence}
A representation $\rho$ is authorization-sufficient on domain $D$ for policy
family $\mathcal P$ when
\[
 \rho(u)=\rho(v) \Longrightarrow
 I_P(\omega(u))=I_P(\omega(v))
\]
for every $u,v\in D$ and $P\in\mathcal P$.
\end{definition}

This definition is policy-relative.  It does not require separating changes
that no admissible policy distinguishes, and it does not privilege a fixed
tuple of atom fields.

\begin{theorem}[Representation collision]
\label{thm:representation-insufficiency}
If $\rho$ is not authorization-sufficient, no deterministic monitor whose only
call-dependent input is $\rho(u)$ can, even when given $P$, both allow every
authorized context and reject every unauthorized context.
\end{theorem}

\paragraph{Proof sketch.}
Insufficiency gives $u,v,P$ with equal representations and opposite ideal
decisions.  The monitor must return the same decision for both.  Allowing
admits the unauthorized member; denying or abstaining withholds the authorized
member.

For any fixed $P$, the number of authorized rows in a representation cell that
also contains an unauthorized row is therefore a lower bound on false denial
or abstention for every sound monitor over that view.  Tool-name, whole-call,
and atom views are all subject to this result; an atom view is useful only when
it preserves the distinctions required by the policy family.

\subsection{Counterfactual Validation}

An intervention $v=do_J(u)$ changes selected call fields, trusted provenance,
or pre-state while holding the tool implementation fixed.  Source execution
observes committed state changes; the trusted harness records policy-relevant
provenance and control metadata.  Together they instantiate $\omega$ on the
tested pair.

\begin{proposition}[Counterfactual witness]
\label{prop:mediation-gap}
If $\rho(u)=\rho(v)$ but some $P\in\mathcal P$ satisfies
$I_P(\omega(u))\ne I_P(\omega(v))$, then $\rho$ is insufficient on every domain
containing $u$ and $v$.
\end{proposition}

The proposition makes counterfactual testing a falsification procedure.  An
output-only change with unchanged $\omega$ is not a witness.  The same committed
state reached under a different trusted provenance may be a witness when an
admissible policy distinguishes that provenance.  Conversely, a representation
change under a true semantic-invariant placebo indicates unnecessary
sensitivity.

For a finite domain, counterexample-guided refinement repeatedly splits a
representation cell using a witnessed separating field or qualifier.  Call a
cell \emph{policy-mixed} when it contains $u,v$ for which at least one
$P\in\mathcal P$ gives different ideal decisions.

\begin{theorem}[Finite refinement termination]
\label{thm:finite-adequacy}
On finite $D$, a refinement procedure that never merges cells can perform at
most $|D|-1$ nonempty cell splits.  If each step separates a witnessed
authorization collision and the procedure stops only when no policy-mixed cell
remains, the resulting representation is authorization-sufficient on
$(D,\mathcal P)$.
\end{theorem}

\paragraph{Proof sketch.}
Every split strictly increases the number of nonempty cells, which is bounded
by $|D|$.  At termination, equal representations imply membership in the same
nonmixed cell.  By the definition of policy-mixed, their ideal decisions agree
for every $P\in\mathcal P$.

Generated interventions need not cover future calls.  We therefore retain
unresolved interventions instead of treating them as evidence of irrelevance.

\subsection{Conditional Atom-Level Mediation}

The representation results become a commit guarantee only when the descriptor,
policy, and enforcement path satisfy separate obligations.  Let $J_P$ be the
ideal policy predicate expressed over a descriptor's atom view.

\begin{theorem}[Conditional atom-level mediation]
\label{thm:conditional-atom-mediation}
Assume four obligations for every reachable effectful call $u$.  First,
descriptor fidelity gives $J_P(A_{d_t}(u))=I_P(\omega(u))$ for every admissible
$P$.  Second, the authorizer returns \Allow{} only when
$J_P(A_{d_t}(u))=1$.  Third, every effectful call is mediated before commit.
Finally, the executor commits exactly the checked call under the same policy
snapshot.  Then every allowed committed call satisfies $I_P(\omega(u))=1$.
\end{theorem}

\paragraph{Proof sketch.}
Consider an allowed committed call $u$.  Complete mediation supplies a check
for $u$, and authorizer soundness gives $J_P(A_{d_t}(u))=1$.  Descriptor
fidelity therefore gives $I_P(\omega(u))=1$.  Call and policy-snapshot
check--use consistency make this conclusion apply at commitment.

The theorem separates representation fidelity from policy correctness.
Counterfactual tests probe the first obligation; ToolSandbox instantiates the
atom interface on a finite domain.  The AgentDojo monitor below enforces a
narrower provenance-origin predicate over registered field values.

\subsection{Implemented Provenance-Origin Confinement}

For a registered call $u$, let $F(u)$ be its security-field values and let
$E(u)$ be its set of registered effect kinds.  Let $U_V(h)$ and $U_E(h)$ be
values and effect requests extracted from provenance-marked untrusted evidence
$h$.  Let $G_V(q)$ and $G_E(q)$ be values and effects grounded in the
authenticated task.  For this monitor, ``introduced solely'' means membership
in the relevant untrusted set without membership in the corresponding grounded
set; it is not a claim of general causal attribution.  The implemented
provenance-origin monitor allows only when
\[
  F(u)\cap U_V(h)\subseteq G_V(q)
  \quad\text{and}\quad
  E(u)\cap U_E(h)\subseteq G_E(q).
\]

\begin{corollary}[Conditional provenance-origin confinement]
\label{thm:single-commit}
Assume that the descriptor covers every effect-bearing field and effect kind
governed by the monitor.  The provenance adapter identifies every untrusted
value and effect request influencing the call, and grounding marks only values
and effects supported by the authenticated task.  Every effectful call is
mediated, and the executor commits exactly the checked call.  Then an allowed
call contains no registered effect or effect-bearing value introduced solely
by untrusted evidence outside the authenticated task.
\end{corollary}

\paragraph{Proof sketch.}
Suppose an allowed call contains such an effect-bearing value.  Descriptor,
provenance, and grounding soundness place it in $F(u)\cap U_V(h)$ but not
$G_V(q)$, contradicting the first allow condition.  If only an effect request
changes, the same obligations place it in $E(u)\cap U_E(h)$ but not $G_E(q)$,
contradicting the second condition.  Mediation and exact execution transfer the
decision to the committed call.

This corollary covers the implemented policy, not ACL authorization,
output-only goals, or dangerous authenticated requests.  Runtime audits test
mediation and check--use consistency; source interventions test descriptor
adequacy on the exercised domain.
Appendix~\ref{app:formal-proofs} gives the complete definitions and proof
details.
